Wind-weighted satellite fire exposure is a common predictor in smoke forecasting, but its
directional component is rarely tested on its own. We build an eleven-year (2015--2025) benchmark
of haze-limited visibility events (visibility below 3\,km with a haze or smoke code) at twelve
airports in Borneo and Sumatra from public METAR archives and VIIRS active-fire detections, and
ask whether wind-directional fire features improve 6\,h and 24\,h forecasts beyond a matched
non-directional exposure. We compare against time-shifted and permuted-wind placebos, controls
with the same missingness as the directional feature, and wind-report categories, under
leave-one-year-out evaluation with station- and region-level block bootstrap. A pre-specified
comparison of single gradient-boosted fits showed a large directional gain at 24\,h
(\pv{loyo24.d.E-D0} average precision), but it was driven by seed-to-seed variation in one
extreme season: single fits of the same model differ by almost 0.1 average precision. With
five-seed ensembles, the directional model is at most a few thousandths of average precision
above its placebos and controls (\pv{ens24.d.E-Rperm1} against the permuted-wind placebo), with no
consistent sign across aggregations, horizons or model families, and no larger with vegetation
fires only. Fire counts improve 24\,h forecasts modestly. We recommend seed
ensembles and mask-preserving placebos for studies that credit wind-aware fire features. Code and
protocol are open.
